\chapter{Hilbert Spaces}

To introduce Hilbert spaces we first recall the definition of a vector space.

\section{Vector space definition and properties}

Definition 4.1. A linear vector space $L$ over a field $F$ (usually either $F=\mathbb{R}$ or $F=\mathbb{C}$) is a set of elements (called vectors):

\begin{equation}
|a\rangle, |b\rangle, |c\rangle, \ldots
\end{equation}

equipped with two operations:
\begin{itemize}
  \item vector addition: $\forall|x\rangle,|y\rangle \in L \longrightarrow |x\rangle+|y\rangle=|z\rangle \in L$
  \item scalar multiplication: $\forall \lambda \in F, \forall|x\rangle \in L \longrightarrow \lambda|x\rangle=|v\rangle \in L$
\end{itemize}

These operations must satisfy the following axioms:
\begin{enumerate}
  \item $|x\rangle+|y\rangle=|y\rangle+|x\rangle$ (commutativity of addition)
  \item $(|x\rangle+|y\rangle)+|z\rangle=|x\rangle+(|y\rangle+|z\rangle)=|x\rangle+|y\rangle+|z\rangle$ (associativity of addition)
  \item $(\lambda\gamma)|x\rangle=\lambda(\gamma|x\rangle)=\lambda\gamma|x\rangle$ (associativity of multiplication)
  \item $(\lambda+\gamma)|x\rangle=\lambda|x\rangle+\gamma|x\rangle, \lambda(|x\rangle+|y\rangle)=\lambda|x\rangle+\lambda|y\rangle$ (distributivity of multiplication)
  \item $|0\rangle+|x\rangle=|x\rangle$ (zero element of addition)
  \item $|x'\rangle+|x\rangle=|0\rangle \Longrightarrow |x'\rangle=-|x\rangle$ (inverse element of addition)
\end{enumerate}

Definition 4.2. A set of vectors is said to be linearly independent if:


\begin{equation}
\sum_{k=1}^{N} c_k|x_k\rangle = 0 \Longleftrightarrow c_k = 0 \;\forall k
\end{equation}

Definition 4.3. A set $\mathcal{B}(L) = \{|x_k\rangle: k \in [1,N]\}$ of $N$ linearly independent vectors forms a basis if any vector $|v\rangle \in L$ can be written as a linear combination of vectors $|x_k\rangle$:

\begin{equation}
|v\rangle = \sum_{k=1}^{N} v_k|x_k\rangle
\end{equation}

Equivalently, a set of linearly independent vectors is said to be complete if it forms a basis.

Definition 4.4. A scalar product (or inner product) is an operation $(x,y)$ that associates a number from the field $F$ to any pair of vectors $|x\rangle$ and $|y\rangle$. From now on, we assume $F = \mathbb{C}$. The scalar product satisfies the following properties:

\begin{enumerate}
  \item $(x, c_1y_1 + c_2y_2) = c_1(x,y_1) + c_2(x,y_2)$
  \item $(c_1x_1 + c_2x_2, y) = c_1^*(x_1,y) + c_2^*(x_2,y)$
\end{enumerate}

Due to the second property (the product is antilinear in the first argument), such a scalar product is said to be sesquilinear. Note that if the field $F = \mathbb{R}$, the scalar product is simply bilinear because $c_i = c_i^*$.

We also define the dual vector:

\begin{equation}
|v\rangle = \sum_{k=1}^{N} v_k|x_k\rangle \longleftrightarrow \langle v| = \sum_{k=1}^{N} v_k^*\langle x_k|
\end{equation}

From this the scalar product can be defined as:

\begin{equation}
\langle a|b\rangle := \sum_{j=1}^{N}\sum_{k=1}^{N} a_j^*b_k\langle x_j|x_k\rangle
\end{equation}

And the norm is defined as:

\begin{equation}
|a| := \sqrt{\langle a|a\rangle} = \sqrt{\sum_{j=1}^{N}\sum_{k=1}^{N} a_j^*a_k\langle x_j|x_k\rangle}
\end{equation}

Definition 4.5. A set of vectors $\mathcal{B}(L)$ forms an orthonormal basis if $\mathcal{B}(L) = \{|k\rangle: k \in [1,N]\}$ is such that:

\[
\langle j|k\rangle = \delta_{jk} = \begin{cases}
0 & \text{if } j \neq k \\
1 & \text{if } j = k
\end{cases}
\]

Thus if we consider an orthonormal vector space the scalar product simplifies to:

\begin{equation}
\sum_{j=1}^{N}\sum_{k=1}^{N} a_j^*b_k\langle x_j|x_k\rangle = \sum_{j=1}^{N}\sum_{k=1}^{N} a_j^*b_k\delta_{jk} = \sum_{k=1}^{N} a_k^*b_k
\end{equation}

And also the norm reduces to:

\begin{equation}
|a| := \sqrt{\langle a|a\rangle} = \sqrt{\sum_{k=1}^{N} a_k^*a_k} = \sqrt{\sum_{k=1}^{N}|a_k|^2}
\end{equation}

To identify the difference between vectors and dual vectors introduced in the Dirac picture of quantum mechanics we can identify them as column vectors and row vectors. In particular the vectors in Dirac's notation also called ket are column vectors, instead the dual vectors also called bra are the row vectors:

\[
|a\rangle = \begin{bmatrix}
a_1 \\
a_2 \\
\ldots \\
a_N
\end{bmatrix} \quad
\langle a| = \begin{bmatrix}
a_1^*, & a_2^*, & \ldots, & a_N^*
\end{bmatrix}
\]

We denote the standard basis of elements as the set of vectors and dual vectors $|k\rangle, \langle k|$ such that:

\[
|k\rangle = \begin{bmatrix}
0 \\
\cdots \\
1 \\
0 \\
\cdots \\
0
\end{bmatrix} \quad
\langle k| = \begin{bmatrix}
0, & \ldots, & 1, & 0, & \ldots, & 0
\end{bmatrix}
\]

Where the only element equal to 1 is the k-th element.

If we choose $F$ as the set of real numbers $\mathbb{R}$ and the basis of vectors:

\begin{align}
&\langle x_1| = (1,0,0) \\
&\langle x_2| = (0,1,0)  \\
&\langle x_3| = (0,0,1)
\end{align}

we are dealing with the Euclidean vector space which is the usual 3D physical space.


\section{Hilbert space}

Let's finally give the definition of a Hilbert space:

Definition 4.6. A Hilbert space is an infinite-dimensional linear vector space over the field $F = \mathbb{C}$ whose vectors have finite norm. Given the standard basis we have:

\begin{equation}
|v\rangle = \sum_{k}^{\infty} v_k|k\rangle
\end{equation}

And so the norm squared must be:

\begin{equation}
|v|^2 = \langle v|v\rangle = \sum_{k=1}^{\infty}|v_k|^2 < \infty
\end{equation}

This means that the series generated from the square modulus of the vector $v$ is the result of a convergent series and so the modulus of the scalar product of any two vectors in a Hilbert space must be finite. This follows from the Cauchy-Schwarz inequality:

\begin{equation}
|\langle a|b\rangle|^2 \leq |a|^2|b|^2 < \infty
\end{equation}

Let us now focus on a particular Hilbert space. We said that all wave functions must satisfy the normalization condition:

\begin{equation}
(\Psi, \Psi) = 1
\end{equation}

And so if we go back to the definition of $L^2$ \eqref{eq:1.50} space we can notice that this space is a vector space that satisfies the properties to be a Hilbert space.

We can notice that the property of additivity of the Hilbert space is expected since the wave functions must satisfy the superposition principle. As we anticipated the space $L^2(V)$ has a scalar product (previously defined). We also have a partially different definition of a complete orthonormal basis for $L^2(V)$:

Definition 4.7. A set of functions $\mathcal{B} = \{\phi_k: k \in [1,\infty]\}$ is a complete (orthonormal) basis of $L^2(V)$ if any wave function is equivalent to its Fourier series:

\begin{equation}
\Psi(\vec{x}) = \sum_{k=1}^{\infty} c_k\phi_k(\vec{x}) \quad \text{where} \quad c_k = (\phi_k, \Psi)
\end{equation}

The formula for the coefficient can be proven using the fact that by construction $(\phi_i, \phi_j) = \delta_{ij}$:

Proof.
\begin{align}
(\phi_k, \Psi) &= \int_V d^3x\,\phi_k^*\Psi = \int_V d^3x\,\phi_k^*\sum_{n=1}^{\infty}c_n\phi_n = \\
&= \sum_{n=1}^{\infty}c_n\int_V d^3x\,\phi_k^*\phi_n = \sum_{n=1}^{\infty}c_n(\phi_k, \phi_n) =  \\
&= \sum_{n=1}^{\infty}c_n\delta_{kn} = c_k
\end{align}

This means that given a wave function and a basis the wave function itself can be identified by an infinite vector containing all the $c_k$ coefficients. Another consequence of the Fourier expansion of the wave function is that the total probability is given by:

\begin{align}
(\Psi, \Psi) &= \int_V d^3x\,\Psi^*\Psi = \int_V d^3x\,\sum_{n=1}^{\infty}c_n^*\phi_n^*\sum_{k=1}^{\infty}c_k\phi_k = \\
&= \sum_{k=1}^{\infty}\sum_{n=1}^{\infty}c_n^*c_k\int_V d^3x\,\phi_n^*\phi_k = \sum_{k=1}^{\infty}\sum_{n=1}^{\infty}c_n^*c_k(\phi_n, \phi_k) = \\
&= \sum_{k=1}^{\infty}\sum_{n=1}^{\infty}c_n^*c_k\delta_{nk} = \sum_{k=1}^{\infty}c_k^*c_k = \sum_{k=1}^{\infty}|c_k|^2  \\
&(\Psi, \Psi) = \sum_{k=1}^{\infty}|c_k|^2
\end{align}

This identity is also known as the Parseval's equation (or Parseval's identity). This equation is true for a generic function only if the base $\mathcal{B}$ is complete, otherwise it can be proven that:

\begin{equation}
(f, f) \geq \sum_{k=1}^{\infty}|f_k|^2
\end{equation}

In the case of wave functions if we take a complete basis we can also state:

\begin{equation}
(\Psi, \Psi) = \sum_{k=1}^{\infty}|c_k|^2 = 1
\end{equation}

And so every term of the series $|c_k|^2$ can be interpreted as an elementary probability related to the eigenstate $\phi_k$ of the wave function. Also, we notice that in general we do not know what the k factor means and it could depend on various physical quantities such as the energy, the angular momentum, etc.

We can now continue our discussion about Hilbert spaces by stating a new theorem:

Theorem 4.1. (Fischer-Riesz) Given a Hilbert space with an orthonormal basis $\mathcal{B} = \{\phi_k \in L^2(V): k \in [1,\infty]\}$ the condition:

\begin{equation}
\sum_{k=1}^{\infty}|c_k|^2 < \infty
\end{equation}

is necessary and sufficient to grant the fact that there exist a function $\Psi \in L^2(V)$ such that:


The completeness relation in Hilbert space provides a powerful framework for understanding quantum systems. When working with an orthonormal basis, we can express any function in the space through its expansion coefficients:

\begin{equation}
\Psi = \sum_{k=1}^{\infty} c_k\phi_k
\end{equation}

\section{Completeness relation}

For a finite-dimensional vector space $L$ with complete basis $\mathcal{B} = \{|k\rangle: k \in [1,N]\}$, the scalar product of any vector with a basis element yields that vector's component along the basis direction:

\begin{equation}
\langle i|v\rangle = \sum_k v_k\langle i|k\rangle = v_i
\end{equation}

This fundamental property allows us to reconstruct any vector through its projections:

\begin{equation}
|v\rangle = \sum_k v_k|k\rangle = \sum_k |k\rangle\langle k|v\rangle = \mathbb{I}|v\rangle
\end{equation}

From this, we derive the completeness relation:

\begin{equation}
\mathbb{I} = \sum_k |k\rangle\langle k|
\end{equation}

The satisfaction of this relation confirms that basis $\mathcal{B}$ is indeed complete. In matrix form, each term $|k\rangle\langle k|$ represents a matrix with a single non-zero element:

\[
|k\rangle\langle k| = \begin{bmatrix}
0 & 0 & \ldots & 0 & 0 \\
0 & \ldots & \ldots & \ldots & 0 \\
\ldots & \ldots & 1 & \ldots & \ldots \\
0 & \ldots & \ldots & \ldots & 0 \\
0 & 0 & \ldots & 0 & 0
\end{bmatrix}
\]

When summed over all values of $k$, these matrices yield the identity:

\[
\mathbb{I} = \begin{bmatrix}
1 & 0 & 0 & 0 & 0 \\
0 & 1 & 0 & 0 & 0 \\
0 & 0 & 1 & 0 & 0 \\
0 & 0 & 0 & 1 & 0 \\
0 & 0 & 0 & 0 & \ldots
\end{bmatrix}
\]

Extending to infinite-dimensional function spaces requires the Dirac delta function. Using the Fourier expansion of a wavefunction:

\begin{align}
\int_V d^3y \left(\sum_k \phi_k^*(\vec{y})\phi_k(\vec{x})\right)\Psi(\vec{y}) &= \sum_k \left(\int_V d^3y\, \phi_k^*(\vec{y})\Psi(\vec{y})\right)\phi_k(\vec{x}) \\
&= \sum_k (\phi_k(\vec{y}), \Psi(\vec{y}))\phi_k(\vec{x}) \\
&= \sum_k c_k\phi_k(\vec{x}) = \Psi(\vec{x})
\end{align}

For this equation to hold, we must have:

\begin{equation}
\sum_k \phi_k^*(\vec{y})\phi_k(\vec{x}) = \delta^3(\vec{y}-\vec{x})
\end{equation}

The Dirac delta function acts on functions in one dimension according to:

\begin{equation}
\int_V dy\, f(y)\delta(y-x) = f(x)
\end{equation}

In three dimensions, this generalizes to:

\begin{align}
\int_{X \times Y \times Z} d^3y\, f(\vec{y})\delta^3(\vec{y}-\vec{x}) &= \int_X dy_1 \int_Y dy_2 \int_Z dy_3\, \delta(y_1-x_1)\delta(y_2-x_2)\delta(y_3-x_3)f(y_1,y_2,y_3) \\
&= f(x_1,x_2,x_3)
\end{align}

Thus, the completeness relation in Hilbert space takes the form:

\begin{equation}
\sum_{k=1}^{\infty} \phi_k^*(\vec{y})\phi_k(\vec{x}) = \delta^3(\vec{y}-\vec{x})
\end{equation}

\section{Examples}
\section{Plane wave basis (finite volume)}
Let's examine a box of volume $V = L^3$ to verify if the basis:


\begin{equation}
\mathcal{B} = \left\{\phi_k = \frac{e^{i\vec{k}\cdot\vec{x}}}{\sqrt{L^3}}: \vec{k} = \frac{2\pi}{L}(n_1, n_2, n_3), n_i \in \mathbb{Z}\right\}
\end{equation}

This set of plane waves exhibits L-periodic boundary conditions, making it particularly useful for analyzing quantum systems in finite volumes. The periodicity manifests as:

\begin{equation}
\phi_k(x_1, x_2, x_3) = \phi_k(x_1+L, x_2, x_3) = \phi_k(x_1, x_2+L, x_3) = \phi_k(x_1, x_2, x_3+L)
\end{equation}

We can verify this property by direct substitution. Taking the first periodicity condition:

\begin{align}
\phi_k(x_1+L, x_2, x_3) &= \frac{e^{i(x_1k_1+Lk_1+x_2k_2+x_3k_3)}}{\sqrt{L^3}} \\
&= \frac{e^{i\vec{k}\cdot\vec{x}}}{\sqrt{L^3}}e^{iLk_1} \\
&= \frac{e^{i\vec{k}\cdot\vec{x}}}{\sqrt{L^3}}e^{i2\pi n_1} = \phi_k(x_1, x_2, x_3)
\end{align}

Since $e^{i2\pi n_1} = 1$ for any integer $n_1$, the periodicity is confirmed.

To establish that these functions form an orthonormal basis, we must verify their orthogonality. The scalar product between two basis elements is:

\begin{align}
(\phi_k, \phi_q) &= \int_V d^3x\, \frac{e^{-i\vec{k}\cdot\vec{x}}}{\sqrt{L^3}} \frac{e^{i\vec{q}\cdot\vec{x}}}{\sqrt{L^3}} \\
&= \int_V d^3x\, \frac{e^{i(\vec{q}-\vec{k})\cdot\vec{x}}}{L^3}
\end{align}

For the case where $\vec{q} = \vec{k}$, this simplifies to:

\begin{equation}
(\phi_k, \phi_q) = \frac{1}{L^3}\int_V d^3x = \frac{V}{V} = 1
\end{equation}

When $\vec{q} \neq \vec{k}$, the integral can be factorized into three one-dimensional integrals:

\begin{equation}
\prod_{j=1}^3 \int_V dx_j\, \frac{e^{i(q_j-k_j)x_j}}{L^3} = \left.\frac{1}{L^3}\prod_{j=1}^3 \frac{e^{i(q_j-k_j)x_j}}{i(q_j-k_j)}\right|_{-L/2}^{L/2}
\end{equation}

Evaluating each factor separately:

\begin{align}
I &= \left.\frac{e^{i(q_j-k_j)x_j}}{i(q_j-k_j)}\right|_{-L/2}^{L/2} \\
&= \frac{e^{i(q_j-k_j)L/2}}{i(q_j-k_j)} - \frac{e^{-i(q_j-k_j)L/2}}{i(q_j-k_j)} \\
&= (\ldots)\sin((q_j-k_j)L/2)  \\
&= (\ldots)\sin((\frac{2\pi}{L}n_j-\frac{2\pi}{L}m_j)L/2) \\
&= (\ldots)\sin(z_j\pi) = 0
\end{align}

Since $z_j$ is an integer, $\sin(z_j\pi) = 0$, confirming orthogonality.

The completeness relation requires:

\begin{align}
\sum_k \phi_k^*(\vec{x})\phi_k(\vec{y}) &= \frac{1}{L^3}\sum_k e^{-i\vec{k}\cdot\vec{x}}e^{i\vec{k}\cdot\vec{y}} \\
&= \frac{1}{L^3}\sum_k e^{i\vec{k}\cdot(\vec{y}-\vec{x})} = \delta^3(\vec{y}-\vec{x})
\end{align}

This identity, a fundamental result in functional analysis, confirms that our basis is complete.
